% Indexed Virtual Number Algebra
% Wisdom Happy
% 2026

\documentclass[12pt]{article}

\usepackage{amsmath, amssymb, amsthm}
\usepackage{hyperref}
\usepackage{booktabs}
\usepackage{enumitem}
\usepackage[margin=1in]{geometry}

% Theorem environments
\newtheorem{theorem}{Theorem}[section]
\newtheorem{lemma}[theorem]{Lemma}
\newtheorem{corollary}[theorem]{Corollary}
\newtheorem{proposition}[theorem]{Proposition}
\theoremstyle{definition}
\newtheorem{definition}[theorem]{Definition}
\newtheorem{axiom}{Axiom}
\newtheorem{example}[theorem]{Example}
\theoremstyle{remark}
\newtheorem{remark}[theorem]{Remark}

% IVNA notation shortcuts
\newcommand{\vz}[1]{0_{#1}}             % indexed zero: \vz{x} -> 0_x
\newcommand{\vi}[1]{\infty_{#1}}        % indexed infinity: \vi{x} -> ∞_x
\newcommand{\vzn}[2]{0^{#1}_{#2}}       % higher-order zero: \vzn{n}{x} -> 0^n_x
\newcommand{\vin}[2]{\infty^{#1}_{#2}}  % higher-order inf: \vin{n}{x} -> ∞^n_x
\newcommand{\ceq}{\mathrel{=\!\!;}}     % collapse equals: \ceq -> =;
\newcommand{\R}{\mathbb{R}}
\newcommand{\C}{\mathbb{C}}

\title{Indexed Virtual Number Algebra: \\
A Structured Interface for Division by Zero \\
and Indeterminate Forms}

\author{Wisdom Happy\thanks{Correspondence: Wisdom@PlayfulSincerity.org. \\[0.5em]
\textit{Methodology:} The core conceptual framework---indexed zeros and infinities, the product rule $\vz{x} \cdot \vi{y} = xy$, and the applications to calculus and indeterminate forms---was conceived and developed by the author beginning in 2017, well before any AI tooling was involved. The formalization, computational verification, and paper preparation were then carried out using the Playful Sincerity Digital Core---an AI-assisted research methodology system built on Claude Code (Anthropic). The Digital Core enforces a Generate--Verify--Revise (GVR) loop---inspired by the Aletheia architecture~\cite{feng2026aletheia}---in which every mathematical claim is verified by at least one independent computational tool and revised if verification fails. The system orchestrates hierarchical planning, parallel agent-based exploration (literature search, symbolic verification, and formal proof running concurrently), and book/paper research integration. Six independent verification tool chains---Python, SymPy, Z3, Lean~4, Wolfram, and a meta-verification layer---produced 403 automated checks with zero failures, each honestly categorized by what it tests (see Appendix~\ref{app:verification}). The complete verification suite is publicly available and independently reproducible at \url{https://github.com/Playful-Sincerity/PS-Research-IVNA}.} \\[0.3em]
Playful Sincerity Research}

\date{March 2026}

\begin{document}

\maketitle

\begin{center}
\textit{``The limit does not exist!''}\\[0.3em]
--- Cady Heron, \textit{Mean Girls} (2004)
\end{center}
\vspace{1em}

\begin{abstract}
We present \textbf{Indexed Virtual Number Algebra (IVNA)}, an algebraic framework
that attaches real-valued indices to zeros and infinities, making division by zero
and indeterminate forms algebraically operable. Where standard mathematics
returns ``undefined'' for expressions like $5/0$, IVNA returns $\vi{5}$---an
indexed infinity that preserves the numerator's identity and permits continued
computation, including exact recovery of the original value via
$\vi{5} \cdot \vz{1} = 5$. More generally, $y / \vz{x} = \vi{\frac{y}{x}}$ for any
real $y$ and index $x$.

The key innovation is the \textbf{indexed product rule}: $\vz{x} \cdot \vi{y} = xy$,
which resolves the product of zero and infinity to a determinate finite number
by tracking provenance through indices. This rule has no exact precedent in
Non-Standard Analysis, grossone, wheel algebra, surreal numbers, or smooth
infinitesimal analysis.

We prove IVNA consistent by constructing an explicit model via embedding in
Robinson's hyperreal numbers, where $\vz{x} = x\varepsilon_0$ for a fixed
infinitesimal $\varepsilon_0$. This embedding is verified computationally
(70 symbolic checks via SymPy, 13 satisfiability and independence checks via Z3,
42 cross-checks via Wolfram) and formally
(core axioms type-checked in Lean~4 with zero proof gaps). IVNA is thus a structured notational
interface to a specific fragment of Non-Standard Analysis---analogous to how
complex number notation ($a + bi$) is an interface to $\R^2$ with a specific
multiplication rule.

We demonstrate applications in limit-free calculus (derivatives of polynomial,
rational, trigonometric, exponential, and logarithmic functions without
$\varepsilon$-$\delta$ arguments), elimination of L'H\^opital's rule,
proportional infinite set sizes, and a ``VEA mode'' for computer arithmetic
that resolves IEEE~754 NaN propagation when operand indices are tracked. The exponential
constant $e$ receives a direct algebraic definition:
$e = (1 + \vz{1})^{\vi{1}}$, revealing the structure of continuous growth
as a scaling symmetry between step size and repetition count.
Beyond calculus, the product rule $\vz{x} \cdot \vi{y} = xy$ recurs as
the resolution mechanism in distribution theory (the Dirac delta),
probability (conditional densities and Bayes' theorem), and algebraic
geometry (a new correspondence between IVNA's index arithmetic and
blow-up resolution)---a structural observation suggesting that these
independently developed techniques share a common algebraic core.
\end{abstract}

% ============================================================
\section{Introduction}
\label{sec:intro}
% ============================================================

\subsection{The Problem}
\label{subsec:problem}

Consider the function $f(x) = 1/x$. As $x$ approaches zero from the right,
$f(x)$ grows without bound. As $x$ approaches zero from the left, $f(x)$
decreases without bound. Standard mathematics concludes that $\lim_{x \to 0} 1/x$
does not exist and that $1/0$ is undefined.

But notice what this conclusion discards. The right-sided and left-sided
behaviors are perfectly well-defined individually---they simply disagree.
The ``undefined'' verdict arises not because nothing is happening, but
because the single symbol $0$ cannot distinguish two different zeros:
the one approached from the right and the one approached from the left.

This is not an isolated case. The expressions $0/0$, $0 \cdot \infty$,
$\infty - \infty$, and $\infty/\infty$ are all classified as
\emph{indeterminate forms}---not because they have no answer, but because
the answer depends on information that the notation has
discarded.\footnote{``Undefined'' is mathematics' way of saying the
question contains more information than the notation can hold.}
When we write $\lim_{x \to 0} \sin(x)/x = 1$ and
$\lim_{x \to 0} x^2/x = 0$, both expressions involve the form $0/0$,
yet they yield different values. The problem is not with zero or infinity.
The problem is that standard notation uses a single symbol for objects
that carry different weights.

The $\varepsilon$-$\delta$ framework of Weierstrass, the hyperreal
construction of Robinson~\cite{robinson1966}, and the grossone methodology
of Sergeyev~\cite{sergeyev2003} each address this problem at different
levels of formality. Yet in none of these frameworks is division by zero
\emph{algebraically operable}---a computation that produces a result you
can continue calculating with and eventually recover the original value from.

In IEEE~754 floating-point arithmetic~\cite{ieee754}, the situation is
worse: \texttt{5.0 / 0.0} returns \texttt{Inf}, and \texttt{Inf * 0.0}
returns \texttt{NaN}---``Not a Number.'' The information that the infinity
came from dividing 5 by zero is destroyed. Every subsequent operation
inherits the \texttt{NaN}, silently corrupting the computation.

\subsection{A Precedent: Complex Numbers}

This situation has a historical parallel. Before the sixteenth century,
$\sqrt{-1}$ was considered impossible---an operation without a legitimate
result. Cardano (1545) encountered it in his cubic formula and called such
quantities ``as refined as they are useless'' (as quoted in~\cite{carlstrom2004}).

Bombelli (1572) took a different approach. He defined consistent rules
for manipulating $\sqrt{-1}$ without knowing what it ``was.'' The rules
worked: imaginary intermediate steps in the cubic formula cancelled out
and produced correct real answers. Gauss (1831) gave the geometric
interpretation---complex numbers as points in a plane---and the objects
became indispensable.

The complex numbers did not introduce new mathematics in the foundational
sense. The ordered pair $(a, b)$ with the multiplication rule
$(a,b)(c,d) = (ac - bd, ad + bc)$ is an operation on $\R^2$. What the
notation $a + bi$ provided was an \emph{interface}: a way to work with
$\R^2$ that made the algebraic structure transparent and the computations
natural. The notation was the contribution.

We propose that division by zero is at a similar stage. The underlying
mathematics exists---Robinson's Non-Standard Analysis provides a rigorous
framework for infinitesimals and infinities. What is missing is an
interface that makes division by zero algebraically operable.

\subsection{IVNA's Proposal}

Indexed Virtual Number Algebra (IVNA) attaches a nonzero \emph{index}
to every zero and every infinity. Instead of a single zero, there is a
family $\{\vz{x} : x \in \C \setminus \{0\}\}$. Instead of a single
infinity, there is a family $\{\vi{x} : x \in \C \setminus \{0\}\}$.
The index $x$ tracks the \emph{weight}---informally, ``how much zero''
or ``how much infinity'' the object represents.

The central rule is the \textbf{indexed product}:
\begin{equation}
\label{eq:product}
\vz{x} \cdot \vi{y} = xy.
\end{equation}
The product of an indexed zero and an indexed infinity is a determinate
finite number---the product of their indices. This resolves the
indeterminate form $0 \cdot \infty$ by preserving the information that
standard notation discards.

Division by zero follows immediately:
\begin{equation}
\label{eq:divzero}
\frac{y}{\vz{x}} = \vi{\frac{y}{x}}.
\end{equation}
The result is an indexed infinity whose index records the numerator and
denominator. Crucially, the operation is \emph{reversible}:
\[
\vi{\frac{y}{x}} \cdot \vz{x} = \frac{y}{x} \cdot x = y.
\]
Multiply the indexed infinity by the indexed zero and you recover
the original value. No information is lost.

The $1/x$ problem from Section~\ref{subsec:problem} resolves cleanly. Approaching zero
from the right means $x = \vz{1}$ (positive index); approaching from
the left means $x = \vz{-1}$ (negative index):
\[
\frac{1}{\vz{1}} = \vi{1}, \qquad
\frac{1}{\vz{-1}} = \vi{-1}.
\]
These are \emph{different indexed infinities}. The two-sided limit ``does
not exist'' because the two sides produce different objects---and IVNA
makes this distinction explicit rather than collapsing both to ``undefined.''

\subsection{Summary of Results}

The contributions of this paper are:

\begin{enumerate}[leftmargin=2em]
\item \textbf{A consistent algebraic framework} (Section~\ref{sec:core})
  in which division by zero produces operable results. Consistency is
  proven via explicit embedding in Robinson's hyperreals
  (Section~\ref{sec:consistency}), verified computationally
  (403 checks across six tools with zero failures),
  and formalized in Lean~4 (11 axioms, 12 theorems).

\item \textbf{A direct algebraic definition of $e$}
  (Section~\ref{sec:calculus}): $e = (1 + \vz{1})^{\vi{1}}$, revealing
  a scaling symmetry in continuous growth.

\item \textbf{Limit-free single-variable calculus}
  (Section~\ref{sec:calculus}): derivatives and integrals computed
  by direct algebraic manipulation, without $\varepsilon$-$\delta$
  arguments or the standard part function. The Fundamental Theorem
  of Calculus reduces to an algebraic identity.

\item \textbf{Elimination of L'H\^opital's rule}
  (Section~\ref{sec:calculus}): the standard indeterminate forms
  $0/0$, $\infty/\infty$, $0 \cdot \infty$, and $\infty - \infty$
  resolved directly by index arithmetic, with $1^\infty$ handled
  by the Exponential Axiom (A-EXP).

\item \textbf{Applications to physics and computer science}
  (Sections~\ref{sec:physics}--\ref{sec:cs}): singularity
  classification in general relativity, and a ``VEA mode'' for
  computer arithmetic that resolves NaN propagation when operand
  indices are tracked.

\item \textbf{Honest positioning} (Section~\ref{sec:literature}):
  IVNA is a structured notational interface to a specific fragment of
  Non-Standard Analysis, analogous to how $a + bi$ is an interface to
  $\R^2$. The indexed product rule $\vz{x} \cdot \vi{y} = xy$ has no
  exact precedent in the reviewed literature, but the underlying
  mathematics is not new.

\item \textbf{Cross-domain structural observation}
  (Section~\ref{sec:crossdomain}): the product rule
  $\vz{x} \cdot \vi{y} = xy$ recurs as the resolution mechanism
  across calculus, distribution theory, probability, and algebraic
  geometry---suggesting that these independently developed
  techniques share a common algebraic core.
\end{enumerate}


% ============================================================
\section{Core Algebra}
\label{sec:core}
% ============================================================

Before proceeding, we note that every axiom and theorem in this section
is machine-verified. The complete Lean~4 formalization, a Python test
suite (30 core tests), and the full verification suite (403 checks
across six tools) are available at:

\begin{center}
\url{https://github.com/Playful-Sincerity/PS-Research-IVNA}
\end{center}

\noindent
To verify independently: clone the repository and run
\texttt{python3 verification/run\_all.py} (requires Python~3.10+,
SymPy, Z3, and optionally Wolfram Engine). This single command runs
all verification categories, checks the Lean~4 build, and saves
detailed per-tool output to \texttt{verification/\_results/}.
We invite skeptical readers to verify before reading further.

\subsection{Virtual Numbers}

We begin with an example. In standard arithmetic, the expression
$5 / 0$ is undefined, and the expression $0 \cdot \infty$ is
indeterminate. In IVNA:
\[
\frac{5}{\vz{1}} = \vi{5}, \qquad
\vi{5} \cdot \vz{1} = 5, \qquad
\vz{3} \cdot \vi{7} = 21.
\]
Division by zero produces an indexed infinity. Multiplying it back
by the indexed zero recovers the original. The product of an indexed
zero and an indexed infinity yields a finite number---the product of
their indices. We now define the objects and rules that make this work.

\begin{definition}[Virtual Numbers]
\label{def:virtual}
A \textbf{virtual number} is either an \emph{indexed zero} $\vz{x}$ or an
\emph{indexed infinity} $\vi{x}$, where $x \in \C \setminus \{0\}$ is the
\textbf{index}. We call $x$ the \emph{weight} of the virtual number.

A \textbf{non-virtual number} is any element of $\R$ (or $\mathbb{C}$).
We write $n$ for a non-virtual number.

The term ``virtual'' reflects the position that exact zeros and exact
infinities do not occur in physical measurement---they are constructs of
mathematical models. This is a naming choice, not a claim about
foundations. The reader may substitute ``indexed'' throughout without
loss of content.
\end{definition}

\begin{remark}[Why ``virtual'']
More nothing will fit into a bigger thing than it would a smaller
thing, even though it is still nothing. A swimming pool drained of
water contains more ``nothing'' than a thimble. The index captures
this: $\vz{5}$ represents the same mathematical zero as $\vz{1}$
under the collapse operator ($\vz{5} \ceq 0 = \vz{1} \ceq 0$),
but it is a zero with five times the capacity. In the NSA embedding
(Section~\ref{sec:consistency}), $\vz{5} = 5\varepsilon_0$ is
literally five times larger than $\vz{1} = \varepsilon_0$---still
infinitesimal, but proportionally larger. These zeros are not the
absence of quantity but the presence of an infinitesimal one.
\end{remark}

\begin{remark}
The index domain is $\C \setminus \{0\}$, not $\C$. The excluded case
$x = 0$ is addressed in Axiom~\ref{ax:indexzero} below: when an
operation produces index zero, the result exits the virtual system
to the real number $0$. This is analogous to $i - i = 0$ exiting
the imaginary numbers. For the core algebra (Sections~\ref{sec:core}--\ref{sec:calculus}),
real indices suffice; complex indices are needed for directional
information (Section~\ref{sec:extended}).
\end{remark}

\subsection{Multiplication}

The multiplication rules define how virtual numbers interact with each
other and with non-virtual numbers. We begin with the central rule.

\begin{axiom}[Zero--Infinity Product]
\label{ax:zi}
$\vz{x} \cdot \vi{y} = xy$
\end{axiom}

This is the core of IVNA. The product of an indexed zero and an
indexed infinity is the product of their indices---a finite,
non-virtual number. By commutativity (formalized in the Lean~4
development), $\vi{y} \cdot \vz{x} = xy$ as well. The indices carry
the information that makes the otherwise indeterminate product
$0 \cdot \infty$ determinate.

\begin{axiom}[Zero--Zero Product]
\label{ax:zz}
$\vz{x} \cdot \vz{y} = \vzn{2}{xy}$
\end{axiom}

\begin{axiom}[Infinity--Infinity Product]
\label{ax:ii}
$\vi{x} \cdot \vi{y} = \vin{2}{xy}$
\end{axiom}

The product of two indexed zeros is a \emph{higher-order} indexed zero
$\vzn{2}{xy}$, and similarly for infinities. The superscript tracks the
order---how many times the virtual number has been ``deepened'' by
multiplication with its own kind. Higher-order virtual numbers interact
with their duals by order cancellation: $\vzn{2}{x} \cdot \vi{y} = \vz{xy}$
(one order cancels), and $\vzn{n}{x} \cdot \vin{n}{y} = xy$ ($n$ orders
cancel completely).

\begin{axiom}[Scalar--Zero Product]
\label{ax:nz}
$n \cdot \vz{x} = \vz{nx}$
\end{axiom}

\begin{axiom}[Scalar--Infinity Product]
\label{ax:ni}
$n \cdot \vi{x} = \vi{nx}$
\end{axiom}

\subsection{Division}

Division rules follow from multiplication by duality. Dividing a
non-virtual number by an indexed zero produces an indexed infinity,
and vice versa. Dividing two virtual numbers of the same kind
produces a non-virtual number---the ratio of their indices.

\begin{axiom}[Division by Zero]
\label{ax:divzero}
$\displaystyle\frac{y}{\vz{x}} = \vi{\frac{y}{x}}$
\end{axiom}

This is the rule that makes division by zero well-defined. The result
is an indexed infinity whose index records both the numerator $y$ and
the denominator's weight $x$. For example, $5 / \vz{1} = \vi{5}$ and
$5 / \vz{2} = \vi{5/2}$.

\begin{axiom}[Division by Infinity]
\label{ax:divinf}
$\displaystyle\frac{y}{\vi{x}} = \vz{\frac{y}{x}}$
\end{axiom}

\begin{axiom}[Zero--Zero Quotient]
\label{ax:zzdiv}
$\displaystyle\frac{\vz{x}}{\vz{y}} = \frac{x}{y}$
\end{axiom}

\begin{axiom}[Infinity--Infinity Quotient]
\label{ax:iidiv}
$\displaystyle\frac{\vi{x}}{\vi{y}} = \frac{x}{y}$
\end{axiom}

The last two axioms resolve two of the classical indeterminate forms.
The expression $0/0$ is indeterminate in standard mathematics because
the single symbol $0$ discards the index information. In IVNA,
$\vz{x}/\vz{y} = x/y$---a determinate ratio. Similarly,
$\vi{x}/\vi{y} = x/y$. The indeterminacy existed because the notation
lacked structure, not because the mathematics lacked answers.

\subsection{Addition and Subtraction}

Virtual numbers of the same kind and order combine by adding indices.
Virtual numbers of different kinds coexist without combining, analogous
to real and imaginary parts in $\mathbb{C}$.

\begin{axiom}[Virtual Addition]
\label{ax:add}
$\vz{x} + \vz{y} = \vz{x+y}$, \quad $\vi{x} + \vi{y} = \vi{x+y}$
\end{axiom}

This is consistent with the scalar rule: $\vz{x} + \vz{y} = \vz{x+y}$
agrees with $1 \cdot \vz{x} + 1 \cdot \vz{y} = \vz{x} + \vz{y}$.

\begin{remark}[Mixed addition]
The sum of a finite number and a virtual number is defined by the
embedding: $n + \vi{x} = \vi{x}$ and $n + \vz{x} = n$ (up to the
collapse operator). A finite quantity is infinitesimal relative to any
indexed infinity and infinite relative to any indexed zero. In the
hyperreal model, this is the standard fact that $r + x\omega_0 \sim
x\omega_0$ for finite~$r$.
\end{remark}

\begin{axiom}[Index Zero Rule (D-INDEX-ZERO)]
\label{ax:indexzero}
When an operation produces index $0$, the result exits the virtual number
system to the real number $0$. The index domain is $\C \setminus \{0\}$.
\end{axiom}

For example, $\vz{3} + \vz{-3} = 0$ (not $\vz{0}$). This is
analogous to $i + (-i) = 0$ exiting the imaginary numbers to the
real line. The rule ensures that when opposing virtual numbers cancel,
the result is the real number zero---not a virtual number with a
degenerate index.

\subsection{Powers and Roots}

\begin{axiom}[Virtual Powers]
$(\vz{x})^n = \vzn{n}{x^n}$, \quad $(\vi{x})^n = \vin{n}{x^n}$
\end{axiom}

\begin{remark}[Notation: order vs.\ index]
The superscript in $\vzn{n}{x}$ denotes the \emph{order}---the number
of infinitesimal factors---not an exponent applied to the index. Thus
$(\vz{x})^2 = \vzn{2}{x^2}$: the order becomes~$2$ and the index
becomes~$x^2$. These are related (raising a virtual number to the
$n$th power produces order~$n$ with index~$x^n$) but conceptually
distinct: the order tracks depth, the index tracks provenance.
\end{remark}

\begin{remark}
Virtual Powers and D-INDEX-ZERO are structural rules that follow
naturally from the algebraic framework. The 11 core axioms formalized
in Lean~4 comprise the nine multiplication and division rules above,
plus the two addition rules (one for indexed zeros, one for indexed
infinities---stated as a single axiom here but formalized separately).
\end{remark}

\subsection{The Collapse Operator}

\begin{definition}[Collapse ($\ceq$)]
The \textbf{collapse operator} $\ceq$ strips indices from virtual numbers,
projecting them to their standard values:
\[
\vz{x} \ceq 0, \qquad \vi{x} \ceq \infty
\]
for all $x \in \C \setminus \{0\}$. This is analogous to the standard part
function $\mathrm{st}()$ in Non-Standard Analysis.
\end{definition}

\subsection{Duality and Reciprocals}

\begin{proposition}[Zero--Infinity Duality]
Indexed zeros and indexed infinities are reciprocals:
\[
\frac{1}{\vz{x}} = \vi{\frac{1}{x}}, \qquad
\frac{1}{\vi{x}} = \vz{\frac{1}{x}}
\]
\end{proposition}

\subsection{Index Domain}

\begin{remark}
The index domain is $\C \setminus \{0\}$. The restriction $x \neq 0$ is
necessary: $\vz{0}$ would represent ``zero with no weight,'' which under
the multiplication rule $\vz{0} \cdot \vi{1} = 0 \cdot 1 = 0$ yields a
result categorically different from $\vzn{2}{1} \cdot \vi{1} = \vz{1}$.
The D-INDEX-ZERO axiom resolves this by routing index-zero results back to
$\R$. Complex indices encode directional information: $0_{e^{i\theta}}$
represents a zero with approach angle $\theta$.

Virtual-valued indices (e.g., $0_{0_z}$) are permitted but reduce to
order changes under the NSA embedding: $0_{0_z} = z\varepsilon_0^2 =
0^2_z$ (a higher-order zero), while $0_{\infty_z} = z$ exits to a real
number. No recursive regress arises---virtual indices are syntactic sugar
for the existing order hierarchy.
\end{remark}


% ============================================================
\section{Extended Axioms}
\label{sec:extended}
% ============================================================

\subsection{The Virtual Taylor Axiom}

\begin{axiom}[Virtual Taylor (A-VT)]
\label{ax:vt}
For any function $f$ analytic at $a$:
\[
f(a + \vz{x}) = f(a) + \vz{f'(a) \cdot x}
  + \vzn{2}{\frac{f''(a) \cdot x^2}{2!}}
  + \vzn{3}{\frac{f'''(a) \cdot x^3}{3!}} + \cdots
\]
Each term carries virtual order $k$ and index $\frac{f^{(k)}(a) \cdot x^k}{k!}$.
Under the collapse operator, all virtual terms vanish: $f(a + \vz{x}) \ceq f(a)$.
\end{axiom}

\subsection{The Exponential Axiom}

\begin{axiom}[Exponential (A-EXP)]
\label{ax:exp}
$(1 + \vz{x})^{\vi{y}} = e^{xy}$
\end{axiom}

\begin{remark}
A-EXP is a corollary of A-VT applied to $f(h) = (1+h)^{1/h}$ in a suitable
sense, but we state it as a separate axiom for clarity. It resolves the
``$e$ problem'': under the basic rules alone,
$(1 + \vz{1})^{\vi{1}} \ceq 1^{\infty} = 1$, which contradicts $e \approx 2.718$.
A-EXP captures the correlation between base perturbation and exponent that the
premature application of $\ceq$ destroys.

The axiom is justified by the NSA embedding (Section~\ref{sec:consistency}):
$\mathrm{st}\!\left((1 + x\varepsilon_0)^{y/\varepsilon_0}\right) = e^{xy}$
is a theorem of hyperreal arithmetic.
\end{remark}

\begin{corollary}
$e = (1 + \vz{1})^{\vi{1}}$.
\end{corollary}

\subsection{Scope}

\begin{remark}[Scope and generalization]
A-VT restricts IVNA's extended operations to analytic functions. Non-analytic
functions (e.g., $|x|$, the Heaviside step function, smooth but non-analytic
functions like $e^{-1/x^2}$) do not extend to virtual arguments via this
mechanism. This covers essentially all functions encountered in single-variable
calculus but excludes distributions, $L^p$ functions, and smooth non-analytic
functions.

The natural generalization is the full transfer principle of
Non-Standard Analysis, which extends \emph{all} first-order
definable functions to the hyperreals. An IVNA-native transfer
axiom---stating that any first-order property of~$\R$ holds for
virtual numbers under the NSA embedding---would remove the
analyticity restriction entirely. We defer this to future work:
the transfer principle requires model-theoretic machinery
(ultrafilters, {\L}o\'s's theorem) that would undermine IVNA's
accessibility goal. The restriction to analytic functions is a
deliberate trade-off between generality and usability.
\end{remark}


% ============================================================
\section{Consistency}
\label{sec:consistency}
% ============================================================

\subsection{The NSA Embedding}

The consistency of IVNA follows from an explicit embedding in Robinson's
hyperreal numbers~\cite{robinson1966}.

\begin{theorem}[IVNA Consistency]
\label{thm:consistency}
IVNA is consistent relative to ZFC with the Axiom of Choice.
\end{theorem}

\begin{proof}[Proof sketch]
Fix a positive infinitesimal $\varepsilon_0$ in the hyperreal field
${}^*\!\R$, and let $\omega_0 = 1/\varepsilon_0$. Define:
\begin{align*}
\vz{x}   &:= x \cdot \varepsilon_0, &
\vi{x}   &:= x \cdot \omega_0 \\
\vzn{n}{x} &:= x \cdot \varepsilon_0^n, &
\vin{n}{x} &:= x \cdot \omega_0^n
\end{align*}
Every IVNA axiom reduces to a tautology of hyperreal arithmetic.
We verify the central rule (Axiom~\ref{ax:zi}):
\[
\vz{x} \cdot \vi{y} = (x\varepsilon_0)(y\omega_0)
= xy \cdot \varepsilon_0 \omega_0 = xy \cdot 1 = xy. \quad\checkmark
\]
The division-by-zero rule (Axiom~\ref{ax:divzero}):
\[
\frac{y}{\vz{x}} = \frac{y}{x\varepsilon_0}
= \frac{y}{x} \cdot \omega_0 = \vi{\frac{y}{x}}. \quad\checkmark
\]
Addition (Axiom~\ref{ax:add}):
\[
\vz{x} + \vz{y} = x\varepsilon_0 + y\varepsilon_0
= (x+y)\varepsilon_0 = \vz{x+y}. \quad\checkmark
\]
The exponential axiom (A-EXP):
\[
(1 + \vz{x})^{\vi{y}} = (1 + x\varepsilon_0)^{y\omega_0}.
\]
Taking logarithms: $y\omega_0 \ln(1 + x\varepsilon_0) =
y\omega_0(x\varepsilon_0 - x^2\varepsilon_0^2/2 + \cdots) = xy +
\text{infinitesimal}$. Applying the standard part:
$\mathrm{st}(e^{xy + \text{inf.}}) = e^{xy}$. $\checkmark$

All remaining axioms are verified analogously.
The full verification is automated across six independent tool
chains---Python, SymPy, Z3, Lean~4, Wolfram, and a meta-verification
layer---totaling 403 checks with zero failures
(Appendix~\ref{app:verification}).
The Lean~4 formalization provides a machine-checked consistency proof.
\end{proof}

\subsection{What IVNA Is}

The embedding reveals IVNA's mathematical identity. The set of all
IVNA expressions maps to the \emph{Laurent monomials} in $\varepsilon_0$:
\[
\{c \cdot \varepsilon_0^k : c \in K \setminus \{0\},\; k \in \mathbb{Z}\},
\]
where $K$ is any field (typically $\C$).
An indexed zero $\vz{x}$ is $x\varepsilon_0$ (a monomial of degree 1).
An indexed infinity $\vi{x}$ is $x\varepsilon_0^{-1}$ (degree $-1$).
A real number $n$ is $n\varepsilon_0^0$ (degree 0). Algebraically,
IVNA is isomorphic to the \emph{unit group} of the Laurent polynomial
ring $K[\varepsilon_0, \varepsilon_0^{-1}]$, which factors as
$K^* \times \mathbb{Z}$: the index $c \in K^*$ carries provenance,
and the grade $k \in \mathbb{Z}$ carries order. Multiplication adds
grades and multiplies indices; the indexed product rule
$\vz{x} \cdot \vi{y} = xy$ is grade-crossing multiplication
$(x, +1) \cdot (y, -1) = (xy, 0)$, exiting to a real number.

This means IVNA is \emph{not} new foundational mathematics. It is a
structured notational interface to a specific, well-understood fragment
of Non-Standard Analysis. The contribution is the interface itself---the
continuously parameterized families of zeros and infinities, the
grade-crossing product rule, and the applications---not the underlying
algebra, which has been studied since Robinson~\cite{robinson1966}.

\subsection{Why This Is Still Valuable}

The complex numbers are ``just'' $\R^2$ with Hamilton's multiplication
rule $(a,b)(c,d) = (ac-bd, ad+bc)$. The ordered-pair construction
preceded the $a + bi$ notation. Yet the notation transformed
mathematics---it made algebraic closure visible, enabled Euler's
formula, and became indispensable in physics and engineering. (We
note a disanalogy: complex multiplication introduced field structure
where $\R^2$ had only vector-space structure, whereas IVNA does not
add algebraic structure beyond what the Laurent monomials already
possess. The analogy is to the \emph{notational interface}, not to
the algebraic enrichment.)

IVNA's contribution is of a similar kind. The indexed product rule
$\vz{x} \cdot \vi{y} = xy$ is not a new theorem. It is a new
\emph{notation} that makes the provenance of zeros and infinities
algebraically transparent. What this notation enables:

\begin{enumerate}
\item Division by zero becomes an operable algebraic step, not an error.
\item The derivative requires no limits, no $\varepsilon$-$\delta$,
  and no standard part function.
\item Indeterminate forms resolve by index arithmetic.
\item Calculators can output $\vi{5}$ instead of \texttt{ERROR}.
\end{enumerate}

The closest prior work is Santangelo's S-Extension~\cite{santangelo2016},
which proposes unique elements per numerator in division by zero.
However, the S-Extension provides no arithmetic for these elements,
no consistency proof, and no applications. IVNA provides all three.


% ============================================================
\section{Applications: Calculus}
\label{sec:calculus}
% ============================================================

\subsection{Polynomial Derivatives}

The IVNA derivative of $f(x) = x^n$ uses the binomial theorem with
virtual perturbation $\vz{1}$:
\[
f'(x) = \frac{f(x + \vz{1}) - f(x)}{\vz{1}}.
\]
For $f(x) = x^2$:
\begin{align*}
f(x + \vz{1}) &= (x + \vz{1})^2 = x^2 + \vz{2x} + \vzn{2}{1} \\
f(x + \vz{1}) - f(x) &= \vz{2x} + \vzn{2}{1} \\
\frac{\vz{2x} + \vzn{2}{1}}{\vz{1}} &= 2x + \vz{1} \ceq 2x.
\end{align*}
No limits are taken. The indexed zero $\vz{1}$ is substituted directly,
the algebra produces the answer, and the collapse operator strips the
remaining virtual term. The computation for $x^n$ proceeds identically
via the binomial theorem: only the $k=1$ term survives after dividing
by $\vz{1}$, yielding $nx^{n-1}$. All higher-order terms produce
virtual zeros that collapse. We have verified this computationally
for $n = 2, 3, 4, 5$ at multiple evaluation points. The verification
exercises the full IVNA pipeline: \texttt{virtual\_taylor()} constructs
the structured virtual sum, Axiom~\ref{ax:zzdiv} divides the leading
term by $\vz{1}$ to extract the derivative, and each residual term is
confirmed to be a higher-order virtual zero that collapses under $\ceq$.

\subsection{Transcendental Derivatives via A-VT}

The Virtual Taylor Axiom (Section~\ref{sec:extended}) extends IVNA
derivatives to all analytic functions. For $f(x) = \sin x$ at $x = a$:
\begin{align*}
\sin(a + \vz{x}) &= \sin a + \vz{\cos(a) \cdot x}
  + \vzn{2}{\tfrac{-\sin(a) \cdot x^2}{2}} + \cdots \\
\frac{\sin(a + \vz{1}) - \sin a}{\vz{1}} &= \cos a + \vz{\cdots}
  \ceq \cos a.
\end{align*}
Similarly: $d/dx(\cos x) = -\sin x$, $d/dx(e^x) = e^x$,
$d/dx(\ln x) = 1/x$, and $d/dx(1/x) = -1/x^2$. Each follows from
substituting $\vz{1}$ into the Taylor expansion at $a$, dividing by
$\vz{1}$, and collapsing. All results are verified computationally.

\subsection{L'H\^opital Elimination}

IVNA makes L'H\^opital's rule unnecessary. Each classical
indeterminate form resolves directly by index arithmetic:

\medskip
\noindent\textbf{The form $0/0$:}
$\displaystyle\lim_{x \to 0} \frac{\sin x}{x}$.
In IVNA: $\sin(\vz{1})/\vz{1}$. By A-VT, $\sin(\vz{1}) = \vz{1}$
(to first order), so $\vz{1}/\vz{1} = 1$. No differentiation needed.

\medskip
\noindent\textbf{The form $\infty/\infty$:}
$\displaystyle\lim_{x \to \infty} \frac{2x+1}{3x+5}$.
In IVNA: $(2\vi{1} + 1)/(3\vi{1} + 5) = \vi{2}/\vi{3} = 2/3$.
The finite terms vanish relative to the indexed infinities (by the
mixed addition remark), and the index ratio gives the answer directly.

\medskip
\noindent\textbf{The form $0 \cdot \infty$:}
$\vz{x} \cdot \vi{y} = xy$ by Axiom~\ref{ax:zi}. The product is
determinate whenever the indices are known.

\medskip
\noindent\textbf{The form $\infty - \infty$:}
$\vi{x} - \vi{y} = \vi{x-y}$. The result is determinate whenever
$x \neq y$, and exits to real $0$ when $x = y$ (D-INDEX-ZERO).

\subsection{Integration}

In IVNA, integration is literal summation. The definite integral
$\int_0^1 f(x)\,dx$ is the sum of $\vi{1}$ terms, each evaluating
$f$ at position $k \cdot \vz{1}$ and weighting by step width $\vz{1}$
(we use $k$ rather than $i$ here to avoid collision with the imaginary
unit used elsewhere in the paper):
\[
\int_0^1 f(x)\,dx = \sum_{k=0}^{\vi{1}} f(k \cdot \vz{1}) \cdot \vz{1}.
\]
For $f(x) = x$:
\begin{align*}
\sum_{k=0}^{\vi{1}} (k \cdot \vz{1}) \cdot \vz{1}
&= \vzn{2}{1} \cdot \sum_{k=0}^{\vi{1}} k
= \vzn{2}{1} \cdot \frac{\vi{1}(\vi{1}+1)}{2} \\
&= \vzn{2}{1} \cdot \frac{\vin{2}{1}}{2}
= \frac{1}{2}. \quad\checkmark
\end{align*}
The key steps: $\vi{1}(\vi{1}+1) = \vin{2}{1}$ because $\vi{1} + 1 = \vi{1}$
(a finite number is negligible relative to an indexed infinity, as justified
by the mixed addition remark in Section~\ref{sec:core}), and
$\vzn{2}{1} \cdot \vin{2}{1} = 1$ by order cancellation.
We have verified $\int_0^1 x^n\,dx = 1/(n+1)$ for $n = 0, 1, 2, 3, 4, 5$.

The Fundamental Theorem of Calculus reduces to an algebraic identity:
adding one term $f(x) \cdot \vz{1}$ to the sum and dividing by $\vz{1}$
recovers $f(x)$. This is not a theorem requiring proof---it is a
tautology of the algebra.

\subsection{The Exponential Constant}

The exponential axiom yields a direct definition:
\[
e = (1 + \vz{1})^{\vi{1}}.
\]
This is not a limit. It is an algebraic expression: one infinitesimal
step ($\vz{1}$), repeated infinitely many times ($\vi{1}$), with
compound growth.

More generally, $(1 + \vz{x})^{\vi{y}} = e^{xy}$ reveals a
\emph{scaling symmetry}: halving the step and doubling the repetitions
gives the same result:
\[
(1 + \vz{x/2})^{\vi{2y}} = e^{(x/2)(2y)} = e^{xy}
= (1 + \vz{x})^{\vi{y}}.
\]
The invariant is the \emph{product} of the indices, not the individual
values. The same exponential $e^c$ can be decomposed infinitely many
ways into step $\times$ repetitions. IVNA makes this decomposition
visible; standard notation hides it.

In physics, where rate $\times$ time $=$ dimensionless exponent,
this becomes: $(1 + \vz{\text{rate}})^{\vi{\text{time}}} =
e^{\text{rate} \times \text{time}}$. The indices carry the physical
dimensions.

Euler's formula extends naturally with complex indices. Setting
$xy = i\pi$ in A-EXP gives Euler's identity:
\[
(1 + \vz{i\pi})^{\vi{1}} = e^{i\pi} = -1.
\]
The most celebrated equation in mathematics is a single axiom
application. The same scaling symmetry from the real case applies:
any factorization of $i\pi$ into step $\times$ repetitions gives
the same result:
\[
(1 + \vz{i})^{\vi{\pi}} = (1 + \vz{\pi})^{\vi{i}}
= (1 + \vz{1})^{\vi{i\pi}} = -1.
\]
More generally, the entire unit circle is parameterized by varying
the imaginary index on the infinity:
\[
(1 + \vz{1})^{\vi{i\theta}} = e^{i\theta}
= \cos\theta + i\sin\theta.
\]
The zero index $i$ gives the step \emph{direction} (imaginary axis);
the infinity index $\theta$ gives the \emph{angle}. Euler's formula
is the statement that an infinitesimal rotation (index~$i$), repeated
infinitely many times (index~$\theta$/step-size), traces a circle. The
indices make the geometric content visible---rotation is iterated
infinitesimal translation, the Lie group/algebra relationship made
algebraically transparent.

\subsection{Residues and Partial Fractions}

The index domain $\C \setminus \{0\}$ unifies IVNA's calculus with
complex analysis. For a rational function $R(z) = P(z)/Q(z)$ with a
simple pole at $z = a$: the zero of $Q$ at $a$ has index $Q'(a)$
(the derivative), so
\[
R(a) = \frac{P(a)}{\vz{Q'(a)}} = \vi{\frac{P(a)}{Q'(a)}}.
\]
The IVNA index $P(a)/Q'(a)$ is exactly the residue $\mathrm{Res}(R, a)$.

\begin{proposition}[Partial Fractions via Index Arithmetic]
\label{prop:partial}
Let $R(z) = P(z)/Q(z)$ be proper with $Q$ having distinct simple
roots $a_1, \ldots, a_n$. The partial fraction decomposition
$R(z) = \sum_{k=1}^n c_k/(z - a_k)$ has coefficients
$c_k = P(a_k)/Q'(a_k)$, each computable by one application of the
IVNA division rule at the corresponding pole. No limit computation,
polynomial long division, or linear system is required.
\end{proposition}

\begin{example}
$R(z) = 1/(z^3 - 1)$. The three poles at $z = 1, \omega, \omega^2$
(cube roots of unity) have IVNA indices
$1/3$, $\omega/3$, $\omega^2/3$. These are the residues. Their sum
$\tfrac{1}{3}(1 + \omega + \omega^2) = 0$ exits to real~$0$ via
D-INDEX-ZERO---recovering the classical fact that residues of a
proper rational function sum to zero.
\end{example}

The same division rule that computes real derivatives
(Section~\ref{sec:calculus}) also computes complex residues and partial
fraction coefficients---a unification made possible by extending the
index domain from $\R$ to $\C$.

\subsection{Proportional Infinite Set Sizes}

In standard set theory, $|\mathbb{N}| = |\mathbb{E}|$ where $\mathbb{E}$
is the even numbers, because a bijection exists. IVNA offers a complementary
perspective: $|\mathbb{N}| = \vi{1}$ and $|\mathbb{E}| = \vi{1/2}$, giving
the ratio $|\mathbb{E}|/|\mathbb{N}| = \vi{1/2}/\vi{1} = 1/2$ by
Axiom~A9 ($\vi{x}/\vi{y} = x/y$).

More generally, for any arithmetic progression with common difference $d$,
the subset has cardinality $\vi{1/d}$, and its ratio to $|\mathbb{N}|$
equals $1/d$---matching the natural density. IVNA does not contradict
Cantor's equipollence; it supplements it with a proportional measure
that distinguishes sets Cantor identifies. This is the same distinction
made by Benci and Di~Nasso's numerosity theory~\cite{benci2003,benci2019},
arrived at through different machinery.


% ============================================================
\section{The Product Rule Across Domains}
\label{sec:crossdomain}
% ============================================================

The indexed product rule $\vz{x} \cdot \vi{y} = xy$
(Axiom~\ref{ax:zi}) was introduced in Section~\ref{sec:core} as an
algebraic rule for resolving $0 \cdot \infty$. In
Section~\ref{sec:calculus}, it appeared as the mechanism behind
derivatives, integration, residue extraction, and compound growth.

A broader pattern emerges: across several independently developed
branches of mathematics, the standard resolution of an indeterminate
form reduces to the same operation---the product (or quotient) of
indexed zeros and indexed infinities. The following table groups
these by genuinely distinct mathematical territory.

\begin{center}
\small
\begin{tabular}{lllll}
\toprule
\textbf{Territory} & \textbf{Instance} & \textbf{Standard form} & \textbf{IVNA form} & \textbf{Ref.} \\
\midrule
\textbf{Calculus}
  & Derivatives       & $\lim \frac{f(x+h)-f(x)}{h}$ & $\frac{\vz{f'(x)}}{\vz{1}}$  & \S\ref{sec:calculus} \\
  & Integration       & $\lim \sum f \Delta x$ & $\sum f \cdot \vz{1}$ over $\vi{1}$ & \S\ref{sec:calculus} \\
  & Compound growth   & $\lim (1+1/n)^n$   & $(1+\vz{1})^{\vi{1}} = e$  & \S\ref{sec:calculus} \\
  & Residues          & $\lim_{z\to a}(z\!-\!a)R(z)$  & $\vz{Q'(a)} \cdot \vi{P(a)/Q'(a)}$  & \S\ref{sec:calculus} \\
  & Removable sing.   & $\lim_{x\to a} f/g$   & $\vz{f'(a)}/\vz{g'(a)}$  & \S\ref{subsec:removable} \\
\midrule
\textbf{Distributions}
  & Dirac delta       & $\lim h \cdot (1/h)$   & $\vz{1} \cdot \vi{1} = 1$  & \S\ref{subsec:dirac} \\
\midrule
\textbf{Probability}
  & Bayes / densities & $f(x,y)/f_X(x)$ & $\vz{f(x,y)}/\vz{f_X(x)}$   & \S\ref{subsec:bayes} \\
\midrule
\textbf{Alg.\ geometry}
  & Blow-up           & Proper transform on $E$ & $(f_a, a)/(g_b, b)$  & \S\ref{subsec:blowup} \\
\bottomrule
\end{tabular}
\end{center}

This is a structural observation, not a claim of mathematical
unification in the sense that group theory unified symmetries or
complex numbers unified $\R^2$ arithmetic. Those unifications were
\emph{generative}---they produced new theorems in their target
domains. IVNA, by its own admission
(Section~\ref{sec:limitations}), proves no new theorems in any
individual domain. What we observe is that a single algebraic
rule---one that falls out of the basic axioms with no additional
machinery---recurs as the computational core of techniques
developed independently across genuinely different branches of
mathematics. The following subsections verify the less familiar
instances.

\subsection{Conditional Densities and Bayes' Theorem}
\label{subsec:bayes}

Let $X, Y$ be continuous random variables with joint density $f_{X,Y}$.
The event $\{X = x\}$ has probability zero---specifically, IVNA
probability $\vz{f_X(x)}$, where $f_X$ is the marginal density.
Bayes' theorem for the conditional density follows from
Axiom~\ref{ax:zzdiv} ($\vz{a}/\vz{b} = a/b$):
\[
f_{Y|X}(y \mid x) = \frac{P(X = x, Y = y)}{P(X = x)}
= \frac{\vz{f_{X,Y}(x,y)}}{\vz{f_X(x)}}
= \frac{f_{X,Y}(x,y)}{f_X(x)}.
\]
The density \emph{is} the index. No additional measure-theoretic
machinery beyond what defines the density itself is required---no
Radon--Nikodym theorem, no limiting argument---just the zero-zero
quotient rule applied to indexed probabilities.

This resolves the Borel--Kolmogorov paradox transparently. The
``paradox'' arises when conditioning on the same geometric event
(e.g., a great circle on $S^2$) using different parameterizations
yields different conditional densities. In IVNA, different
parameterizations produce different indexed zeros---$\vz{f_X(x)}$
vs.\ $\vz{g_U(u)}$---and Axiom~\ref{ax:zzdiv} correctly produces
different quotients. The result is not paradoxical; it is the expected
consequence of dividing different indexed zeros.

Verification: tested on bivariate normal ($\rho$ general),
Gumbel bivariate exponential ($\theta = 0.5$), and bivariate
Cauchy (no finite moments). All conditionals integrate to~1.
The Cauchy case is the strongest test: a distribution with no
mean or variance still produces correct conditional densities
via Axiom~\ref{ax:zzdiv} (6 checks, 0 failures; see supplementary
\texttt{verify-01-bayes-\allowbreak{}theorem.md}).

The closest prior work is Jacobs~\cite{jacobs2021}, who uses
infinitesimal ratios for conditional densities within a
categorical probability framework. IVNA's contribution is the
directness: Axiom~\ref{ax:zzdiv} \emph{is} Bayes' theorem, stated
as a single axiom that was already present in the core algebra for
independent reasons.

\subsection{The Dirac Delta from the Product Rule}
\label{subsec:dirac}

The Dirac delta ``function'' $\delta(x)$ satisfies $\int \delta(x)\,dx = 1$
and $\delta(x) = 0$ for $x \neq 0$. Distribution theory
(Schwartz, 1950) formalizes this as a linear functional, not a function.

In IVNA, $\delta$ is a function: at $x = 0$, it takes the value $\vi{1}$
(infinite height), defined on a domain of width $\vz{1}$ (infinitesimal
support). Its ``area'' is:
\[
\vi{1} \cdot \vz{1} = 1 \quad \text{(Axiom~\ref{ax:zi})}.
\]
The four standard properties follow algebraically:
\begin{enumerate}
\item \textbf{Normalization}: $\vi{1} \cdot \vz{1} = 1$. \checkmark
\item \textbf{Sifting}: $f(0) \cdot \vi{1} \cdot \vz{1} = f(0)$. \checkmark
\item \textbf{Scaling}: $\delta(ax)$ has height $\vi{1/a}$, width $\vz{a}$,
  so $\vi{1/a} \cdot \vz{a} = 1/|a| \cdot a/a = 1/|a|$. \checkmark
\item \textbf{Convolution}: index products of two deltas compose. \checkmark
\end{enumerate}

For any nascent delta family with height $h(\varepsilon)$ and width
$w(\varepsilon)$, the product $h \cdot w = 1$ holds at every
$\varepsilon > 0$---not just in the limit. Axiom~\ref{ax:zi}
characterizes the \emph{entire family}, not its limit.

Verification: 28 checks across rectangular, Gaussian, and Lorentzian
nascent delta families (see supplementary
\texttt{verify-03-dirac-delta.md}).

The NSA treatment of delta functions dates to
Robinson~\cite{robinson1966}; Todorov~\cite{todorov1990} and
Vernaeve~\cite{vernaeve2025} develop this further. IVNA's
contribution is that the single axiom~\ref{ax:zi} unifies all four
properties, and the characterization is exact equality (not
the ``infinitely close'' relation $\approx$ of NSA).

\subsection{Removable Singularities as Index Cancellation}
\label{subsec:removable}

A removable singularity of $f(x)/g(x)$ at $x = a$ occurs when both
$f$ and $g$ vanish at $a$ with the same order. In IVNA:
\[
\frac{f(a)}{g(a)} = \frac{\vz{f'(a)}}{\vz{g'(a)}} = \frac{f'(a)}{g'(a)}.
\]
The singularity ``removes itself'' by index cancellation
(Axiom~\ref{ax:zzdiv}). No limit is computed; the value at the
point is defined by the algebra.

More generally, if $f$ vanishes to order $m$ and $g$ to order $n$ with
$m = n$, the IVNA quotient $\vzn{m}{c_f}/\vzn{n}{c_g} = c_f/c_g$
exits to a finite number. When $m > n$, the result is $\vz{c_f/c_g}$
(the function has a zero). When $m < n$, the result is $\vi{c_f/c_g}$
(a genuine pole). The order comparison determines the singularity type;
the index carries the value.

Verification: 24 checks across polynomial and trigonometric examples
(see supplementary \texttt{verify-04-removable-singularities.md}).

\subsection{The Blow-Up Correspondence}
\label{subsec:blowup}

In algebraic geometry, \emph{blow-ups} resolve indeterminacy at singular
points by replacing a point with the projective space of approach
directions (the exceptional divisor $E \cong \mathbb{P}^1$). The
following theorem establishes a precise correspondence.

\begin{theorem}[IVNA--Blow-Up Correspondence]
\label{thm:blowup}
Let $f, g$ be polynomials vanishing at the origin with orders $a, b$
and leading homogeneous forms $f_a, g_b$. In the blow-up chart $y = tx$,
the proper transform of $f/g$ restricts to $E$ as the rational function
$f_a(1,t)/g_b(1,t)$. The IVNA quotient $\vzn{a}{f_a}/\vzn{b}{g_b}$
produces the same data: index $f_a/g_b$ and grade $a - b$.
\end{theorem}

\begin{proof}
\emph{Blow-up side.} Substituting $y = tx$ into $f$ factors out $x^a$;
into $g$ factors out $x^b$. The proper transform $x^{b-a} f_a(1,t)/g_b(1,t)$
restricts to $E$ (where $x = 0$) as $f_a(1,t)/g_b(1,t)$ when $a = b$,
and vanishes or diverges along $E$ when $a \neq b$.

\emph{IVNA side.} Under the $K^* \times \mathbb{Z}$ characterization,
$(f_a, a) / (g_b, b) = (f_a/g_b,\, a - b)$. The grade $a - b$
determines the virtual type (zero if positive, real if zero, infinity
if negative); the index $f_a/g_b$ carries the resolved value.

The $\mathbb{P}^1$ coordinate $t = y/x$ parametrizes both: the IVNA
index at direction $t$ equals the blow-up proper transform evaluated at
the corresponding point on $E$.
\end{proof}

\begin{example}
For $f = xy$, $g = x^2 + y^2$ (both order 2): blow-up gives
$t/(1+t^2)$ on $E$; IVNA gives $\vz{xy}/\vz{x^2+y^2} = xy/(x^2+y^2)$.
At direction $t$, both evaluate to $t/(1+t^2)$.
\end{example}

The key distinction is structural: a blow-up produces a geometric
object (the exceptional divisor) within scheme-theoretic framework;
IVNA produces an algebraic element supporting further multiplication,
division, and exponentiation. Blow-ups resolve singularities globally;
IVNA provides arithmetic on singular values locally. Neither subsumes
the other. Notably, IVNA's purely algebraic axioms apply over any
field---including positive characteristic, where resolution of
singularities in dimension $\geq 4$ remains open.

No prior division-by-zero framework (meadows, wheels, transreal
arithmetic, S-Extensions) has connected to blow-up theory; the
correspondence appears to be new. Verification: 14 checks across
polynomial and trigonometric examples (see supplementary materials).

\subsection{A Structural Remark}

The observation that Axiom~\ref{ax:zi} ($\vz{x} \cdot \vi{y} = xy$)
independently resolves indeterminate forms across nine domains can be
understood algebraically: each domain's resolution mechanism is an
independent embedding into IVNA's graded algebra $K^* \times \mathbb{Z}$,
where the product rule serves as the universal resolution step. The
index carries domain-specific information (derivative coefficients,
probability densities, residues, blow-up coordinates), while the
grade crossing from virtual to real ($+1 + (-1) = 0$) performs the
resolution.

A natural objection: $K^* \times \mathbb{Z}$ is one of the simplest
algebraic structures, and \emph{any} domain involving quantities that
vanish or diverge will have Laurent-like behavior at singularities.
The cross-domain pattern may therefore reflect the structure of
singularities in general rather than a deep connection between the
specific domains listed. We acknowledge this. The observation's value
is not that the pattern is surprising in retrospect, but that it
was not previously articulated: no prior work collects these domain
resolutions under a single named algebraic operation. Whether this
articulation has deeper structural significance---a functorial
relationship between domains, perhaps---is an open question for
future work.


% ============================================================
\section{Applications: Physics}
\label{sec:physics}
% ============================================================

IVNA's physics applications are primarily notational---they provide
cleaner ways to express and classify singularities, not new physical
predictions. We state this explicitly to avoid overclaiming.

\subsection{Singularity Classification}

Physical singularities arise when a quantity diverges. In IVNA, the
\emph{order} of the indexed infinity encodes the divergence rate,
and the \emph{index} encodes the physical parameters:

\begin{center}
\begin{tabular}{lll}
\toprule
\textbf{Singularity} & \textbf{Standard} & \textbf{IVNA} \\
\midrule
Coulomb ($r \to 0$) & $F \to \infty$ & $F = \vin{2}{kq_1q_2/x^2}$ \\
Newton gravity & $\phi \to -\infty$ & $\phi = -\vi{GM/x}$ \\
Schwarzschild ($r \to 0$) & $K \to \infty$ & $K = \vin{6}{\frac{48G^2M^2}{x^6}}$ \\
Point charge self-energy & $U \to \infty$ & $U = \vi{kq^2/x}$ \\
\bottomrule
\end{tabular}
\end{center}

The IVNA order distinguishes singularity severity: Schwarzschild
(order~6) is ``more singular'' than Coulomb (order~2). The index
preserves the physical parameters---charge, mass, coupling constants---that
standard ``$\to \infty$'' discards.

Ratios of singularities yield finite, physically meaningful quantities.
For two Coulomb forces at the same point:
$\vin{2}{kq_1q_2/x^2} / \vin{2}{kq_3q_4/x^2} = q_1q_2 / q_3q_4$.

\subsection{What IVNA Does Not Do for Physics}

IVNA cannot replace regularization in quantum field theory. The
subtraction $\vi{a} - \vi{b} = \vi{a-b}$ is too simple to capture
nested divergences, running couplings, or the combinatorial structure
of Feynman diagrams. Renormalization requires dimensional regularization,
Pauli--Villars, or similar techniques that IVNA does not provide.

IVNA also cannot address topological singularities (cosmic strings,
conical deficits) where the divergence is geometric rather than
algebraic.


% ============================================================
\section{Applications: Computer Science}
\label{sec:cs}
% ============================================================

\subsection{VEA Mode}

In IEEE~754 floating-point arithmetic~\cite{ieee754}, division by zero
produces \texttt{Inf}, and \texttt{Inf * 0} produces \texttt{NaN}. Once
a \texttt{NaN} appears, it propagates through every subsequent operation,
silently corrupting the computation.

IVNA suggests an alternative we call \textbf{VEA mode} (Virtual Extended
Arithmetic). In VEA mode, \texttt{5.0 / 0.0} returns $\vi{5}$---an
indexed infinity that remembers its origin. Multiplying $\vi{5}$ by
$\vz{1}$ recovers $5$. No information is lost.

\begin{center}
\begin{tabular}{lll}
\toprule
\textbf{Expression} & \textbf{IEEE 754} & \textbf{VEA(IVNA)} \\
\midrule
\texttt{5 / 0}      & \texttt{Inf}  & $\vi{5}$ \\
\texttt{-3 / 0}     & \texttt{-Inf} & $\vi{-3}$ \\
\texttt{Inf * 0}    & \texttt{NaN}  & $5$ (if $\vi{5} \cdot \vz{1}$) \\
\texttt{Inf - Inf}  & \texttt{NaN}  & $\vi{a-b}$ (indices known) \\
\texttt{0 / 0}      & \texttt{NaN}  & $x/y$ (indices known) \\
\bottomrule
\end{tabular}
\end{center}

\begin{proposition}[NaN Elimination]
\label{prop:nan}
Under VEA mode, every IEEE~754 NaN-producing primitive binary operation
between well-indexed virtual numbers has a determinate result:
$\vz{x}/\vz{y} = x/y$,\quad
$\vz{x} \cdot \vi{y} = xy$,\quad
$\vi{x}/\vi{y} = x/y$,\quad
$\vi{x} - \vi{y} = \vi{x-y}$ (or $0$ when $x = y$).
\end{proposition}

\begin{proof}
Each case is a direct application of Axioms~A3, A8, A9, and A10--A11
respectively. All eight NaN-producing primitive operations are
computationally verified with zero failures.
\end{proof}

\begin{example}[Quadratic formula]
Solving $x^2 - 2 \times 10^8 x + 1 = 0$: the standard formula
produces $x_2 = 10^8 - \sqrt{10^{16} - 1}$, which in float64
returns exactly $0$ (100\% wrong; the true value is $5 \times 10^{-9}$).
In VEA mode, the near-cancellation produces an indexed zero
$\vz{10^{-8}}$, preserving the residual's identity through subsequent
operations and recovering the exact root.
\end{example}

A working VEA calculator prototype is included in the supplementary
repository.

\medskip
\noindent\textbf{Comparison with existing tools.}

\begin{center}
\begin{tabular}{lcccc}
\toprule
& \textbf{Determ.\ $0/0$} & \textbf{Past NaN} & \textbf{Roundtrip} \\
\midrule
IEEE 754              & --  & --  & -- \\
Interval arithmetic   & --  & P   & -- \\
Significance arith.   & D   & --  & -- \\
Dual numbers / AD     & P   & Y   & -- \\
TwoSum / EFTs         & --  & --  & P \\
Wheel algebra         & Y$^\ddagger$   & --  & -- \\
\textbf{IVNA (VEA)}   & \textbf{Y} & \textbf{Y} & \textbf{Y} \\
\bottomrule
\end{tabular}
\end{center}

{\small Y = yes, P = partial, D = diagnoses only, -- = no.
$^\ddagger$Wheel algebra gives $\bot$ for $0 \cdot \infty$
(absorbing, not determinate).}

\medskip
Fog~\cite{fog2019} has proposed encoding exception metadata in NaN
payloads for the ForwardCom architecture. VEA differs: Fog's payloads
track \emph{where} the exception occurred (code address), while IVNA's
indices track \emph{what} the mathematical provenance is (numerator,
denominator). VEA indices are algebraically operable; Fog's payloads
are not.


% ============================================================
\section{Literature Positioning}
\label{sec:literature}
% ============================================================

\subsection{Related Frameworks}

Six established frameworks address overlapping concerns:

\begin{enumerate}
\item \textbf{Non-Standard Analysis} (Robinson, 1966)~\cite{robinson1966}:
  infinitesimals via hyperreals. Keisler~\cite{keisler1976} developed
  the pedagogical calculus curriculum. Division by zero remains undefined;
  $0 \times \infty$ is indeterminate.

\item \textbf{Grossone}
  (Sergeyev, 2003)~\cite{sergeyev2003,sergeyev2017,sergeyev2019}:
  a numeral for $|\mathbb{N}|$, with its reciprocal as infinitesimal.
  No division by zero; no general $0 \times \infty$ resolution.
  Subject to criticism~\cite{gutman2017}.

\item \textbf{Numerosity} (Benci--Di Nasso, 2003)~\cite{benci2003,benci2019}:
  proportional sizes of infinite sets. No division by zero or calculus.

\item \textbf{Wheel algebra} (Carlstr\"om, 2004)~\cite{carlstrom2004}:
  defines $1/0 = \infty$, but $0 \cdot \infty = \bot$ (absorbing
  element). Information is destroyed, not preserved.

\item \textbf{Surreal numbers} (Conway, 1976)~\cite{conway1976}:
  the largest ordered field. Division by zero undefined.

\item \textbf{SIA/SDG} (Kock, Bell)~\cite{kock2006,bell1998}:
  nilsquare $\varepsilon$ with $\varepsilon^2 = 0$. Requires
  intuitionistic logic.
\end{enumerate}

Additional approaches include meadows~\cite{bergstra2009}, which
formalize division-by-zero via equational specifications;
Saitoh's generalized division framework~\cite{saitoh2014}; and
Bergstra's survey~\cite{bergstra2019} of the design space.
Wenmackers~\cite{wenmackers2024} provides a philosophical comparison
of methods for comparing infinite quantities.

The IVNA--Blow-Up Correspondence (Theorem~\ref{thm:blowup} in
Section~\ref{subsec:blowup}) establishes a precise connection
between IVNA's index arithmetic and blow-up resolution in algebraic
geometry. No prior division-by-zero framework has made this
connection.

\subsection{Comparison}

\begin{center}
\begin{tabular}{lccccccc}
\toprule
& NSA & Gr. & Num. & Wheel & Sur. & SIA & \textbf{IVNA} \\
\midrule
Div.\ by zero    & -- & -- & -- & P & -- & -- & \textbf{Y} \\
$0 \times \infty$ resolved & -- & P & -- & -- & -- & -- & \textbf{Y} \\
Limit-free deriv. & Y$^*$ & P & -- & -- & P & Y$^\dagger$ & \textbf{Y} \\
Prop.\ set sizes & -- & Y & Y & -- & -- & -- & \textbf{Y} \\
Classical logic  & Y & Y & Y & Y & Y & -- & \textbf{Y} \\
\bottomrule
\end{tabular}
\end{center}

{\small Y = yes, P = partial, -- = no.
$^*$Requires standard part function.
$^\dagger$Requires intuitionistic logic.}

\medskip
No single existing framework provides all five capabilities.

\subsection{Genuine Novelty}

The indexed product rule $\vz{x} \cdot \vi{y} = xy$ has no exact
precedent. In NSA, $\varepsilon \cdot \omega$ is indeterminate without
specifying which pair. In wheel algebra, $0 \cdot \infty = \bot$.
IVNA's rule is a general algebraic law over a continuously parameterized
index space. A systematic search of arXiv, Semantic Scholar, Google
Scholar, and the general web (approximately 125 targeted queries across
six claim streams) confirmed that the terms ``indexed zeros'' and
``indexed infinities'' together return zero results in the academic
literature.

Of the new results presented in this paper, two appear to be fully
novel: the IVNA--Blow-Up Correspondence (Theorem~\ref{thm:blowup})
and the cross-domain structural observation
(Section~\ref{sec:crossdomain}). No prior division-by-zero framework
has connected to blow-up theory, and no prior work has identified the
same algebraic operation recurring across nine domains.

Four results are partially anticipated by prior work but with
meaningful IVNA-specific differentiation:
the probability/Bayes connection (cf.\ Jacobs~\cite{jacobs2021}),
the Dirac delta treatment (cf.\ Todorov~\cite{todorov1990},
Vernaeve~\cite{vernaeve2025}),
the infinity subtraction / renormalization connection
(cf.\ Albeverio et al.~\cite{albeverio1986}),
and the $K^* \times \mathbb{Z}$ algebraic characterization
(cf.\ Santangelo~\cite{santangelo2016}).
In each case, IVNA's contribution is the directness and algebraic
operability of the result, not the underlying mathematical fact.

The closest prior work remains Santangelo's
S-Extension~\cite{santangelo2016}, which proposes unique elements per
numerator in division by zero. The S-Extension establishes the
structural intuition but provides no arithmetic, no consistency proof,
and no applications. IVNA provides all three.
Fereydoni~\cite{fereydoni2025} independently introduces stratified
infinities $\{\infty_k\}$ via fixed-point iteration on a compactified
metric space, with notation superficially similar to IVNA's
$\vi{x}$; however, Fereydoni's indices encode algebraic hierarchy
(iterated fixed points), not the zero or blow-up that produced the
infinity, and no product rule or cross-domain applications are
developed.

\subsection{Addressing Potential Criticisms}

We anticipate five objections:

\begin{enumerate}
\item \emph{``Just notation for NSA.''} Correct. And $a + bi$ is just
  notation for $\R^2$. The notation is the contribution.

\item \emph{``Santangelo (2016) did this.''} Santangelo proposed the
  structure; IVNA provides the complete package.

\item \emph{``Axiom of Choice dependency.''} Yes, via the ultrafilter
  for the NSA embedding. Standard in mainstream mathematics.

\item \emph{``Grossone-type fringe work.''} Five differences:
  no circularity, no independence claim, decidable comparisons,
  working implementation, Lean~4 proofs.

\item \emph{``A-VT is ad hoc.''} It is the IVNA translation of Taylor
  expansion in ${}^*\!\R$. Scope explicitly restricted to analytic
  functions; generalization via the transfer principle is identified
  as future work (Section~\ref{sec:extended}).
\end{enumerate}


% ============================================================
\section{Limitations and Scope}
\label{sec:limitations}
% ============================================================

\subsection{What IVNA Does Not Do}

\begin{enumerate}
\item \textbf{Non-analytic functions.} A-VT requires analyticity.
  Functions like $|x|$ do not extend to virtual arguments.

\item \textbf{New theorems.} Every calculus result here is known from
  standard analysis or NSA. IVNA proves nothing new---it proves known
  things more simply.

\item \textbf{Nonlinear ODEs.} Linear ODEs telescope cleanly via
  A-EXP. Nonlinear ODEs reduce to forward Euler---correct but not
  insightful.

\item \textbf{Renormalization.} IVNA's singularity notation is too
  simple for QFT's combinatorial divergences.

\item \textbf{Foundational claims.} None. IVNA is consistent relative
  to ZFC+AC, not a foundation for mathematics.

\item \textbf{Sub-first-order infinities.} IVNA's indexed infinities
  $\vi{x} = x/\varepsilon_0$ are first-order (linear in
  $1/\varepsilon_0$). The harmonic series $H_n \sim \ln n + \gamma$
  diverges to an infinity of \emph{lower} order than any $\vi{x}$:
  $\ln(1/\varepsilon_0)$ grows slower than $x/\varepsilon_0$ for any
  $x > 0$. IVNA's current notation does not parameterize logarithmic
  or sub-logarithmic infinities. This precisely delineates the
  scope of the indexed notation.

\item \textbf{Restatements vs.\ new results.} Several results in
  Section~\ref{sec:crossdomain} are IVNA restatements of facts
  known from standard analysis or NSA. The Dirac delta properties
  follow from NSA (Robinson, 1966); the conditional density
  formula is Bayes' theorem in different notation. What IVNA adds
  in each case is directness and algebraic operability, not the
  underlying mathematical fact. We state this explicitly because
  honest positioning is more valuable than overclaiming.
\end{enumerate}

\subsection{Is IVNA ``Just Notation''?}

IVNA is isomorphic to Laurent monomials in a reference infinitesimal.
Is it ``just notation''?

Complex number notation is ``just notation'' for $\R^2$. It did not add
new theorems to real analysis. It added a way of thinking that made
algebraic closure visible, enabled Euler's formula, and became the
language of quantum mechanics. The contribution was the interface.

IVNA's contribution is of the same kind. The indexed product rule
$\vz{x} \cdot \vi{y} = xy$ is a theorem of hyperreal arithmetic.
What IVNA adds is a notation that makes this theorem---and its
consequences---accessible to anyone who can multiply.


% ============================================================
\section{Research Methodology}
\label{sec:methodology}
% ============================================================

The development of IVNA employed a structured AI-assisted research
methodology that may be of independent interest.

\textbf{Generate--Verify--Revise (GVR) loop.} Every mathematical
claim was first generated from theoretical reasoning, then verified
by at least two independent computational tools (drawn from Python,
SymPy, Z3, Lean~4, and Wolfram Mathematica), and revised if
verification failed. This protocol, inspired by the Aletheia
architecture~\cite{feng2026aletheia}, enforces falsifiability at the
point of creation rather than after publication.

\textbf{Adversarial debate.} Key thesis claims were stress-tested
via structured adversarial analysis: independent AI agents assigned
PRO and CON positions argued from the paper's source text, with a
neutral agent synthesizing. Across four debates, this process was
not merely verificatory but \emph{generative}. The first debate
(March 2026) identified blow-up theory as a comparator that IVNA
had not engaged; the resulting research produced the IVNA--Blow-Up
Correspondence (Theorem~\ref{thm:blowup}), the paper's most novel
result. A subsequent debate discovered that the cross-domain
observation existed only in supplementary research documents and
had not been incorporated into the paper, directly precipitating
Section~\ref{sec:crossdomain}. A final debate calibrated the
framing from ``unification'' to ``structural observation''---a
distinction this section's careful hedging reflects. Each debate
exposed a specific gap; closing that gap produced the next real
finding.

\textbf{Multi-tool cross-verification.} High-priority claims (the
Bayes/Axiom~\ref{ax:zzdiv} identification, the Borel--Kolmogorov
dissolution, the Dirac delta properties) were verified with three
independent tool chains rather than two. Tool disagreement triggers
investigation, not averaging.

\textbf{Meta-verification.} A post-completion audit revealed that the
original verification suite conflated IVNA-native checks (exercising the
\texttt{Virtual} class) with classical computations narrated in IVNA
notation. The suite was rebuilt with three explicit categories:
Category~A (IVNA-native, using the axiom implementation), Category~B
(NSA embedding consistency via SymPy), and Category~C (classical
correspondence, honestly labeled). A meta-verification layer checks
that the suite itself is structurally sound---catching tautological
assertions, circular passthrough functions, and category violations.
This ``meta-GVR'' step---verifying that verification tests what it
claims to test---is itself a methodological finding: automated
verification can produce false confidence when tests exercise notation
rather than the target algebra.

The complete methodology, including all agent transcripts, debate
records, and verification logs, is available in the supplementary
repository. A detailed case study of the AI-assisted research
process is in preparation as a companion document.


% ============================================================
\section{Conclusion}
\label{sec:conclusion}
% ============================================================

\subsection{Summary}

Indexed Virtual Number Algebra is a consistent, computationally verified,
and formally proven algebraic framework for zeros and infinities. Its
central rule---$\vz{x} \cdot \vi{y} = xy$---resolves indeterminate forms
by preserving index information through operations. Consistency is proven
via NSA embedding, with 403 automated checks across six independent
tool chains and Lean~4 formalization, with zero failures.

The same product rule that defines IVNA's core algebra recurs as the
resolution mechanism across calculus, distribution theory, probability,
and algebraic geometry (Section~\ref{sec:crossdomain}). This structural
observation---verified but not yet fully explained---suggests that the
indexed product rule captures something about how mathematics resolves
singularities, though whether this reflects deep structure or the
inherent simplicity of $K^* \times \mathbb{Z}$ remains an open
question.

IVNA is not new foundational mathematics. It is a structured interface to
Non-Standard Analysis, analogous to $a + bi$ for $\R^2$. Its value is in
making existing mathematics more accessible and computable at the boundary
between zero and infinity.

\subsection{Future Work}

\begin{enumerate}
\item \textbf{Integration formalization} in Lean~4.
\item \textbf{VEA implementation} in a production language with IEEE~754
  benchmarks.
\item \textbf{Complex indices}: systematic treatment of directional
  singularities where the index encodes approach direction.
\item \textbf{Pedagogical testing}: controlled comparison of IVNA-based
  vs.\ standard calculus instruction.
\item \textbf{Transfer principle generalization}: replace the analyticity
  restriction of A-VT with a full IVNA-native transfer axiom, extending
  virtual arguments to all first-order definable functions. This requires
  engaging with model-theoretic foundations (ultrafilters, \L o\'s's theorem)
  while preserving IVNA's notational accessibility.
\item \textbf{Reparametrization invariance}: if $f(x) = a \cdot h(x)$
  and $g(x) = b \cdot h(x)$ for any $h$ with $h(c) = 0$, the IVNA index
  of $f/g$ at $c$ reduces to $a/b$ regardless of the shared factor. This
  algebraic invariance property---defined at the point, not via limits---is
  a candidate for a formal theorem connecting IVNA's index algebra to
  classical singularity classification.
\item \textbf{Measure-theoretic formalization}: establish the
  connection between IVNA's indexed zero probabilities and the
  Radon--Nikodym derivative. Determine whether the
  Borel--Kolmogorov dissolution
  (Section~\ref{subsec:bayes}) holds under weaker regularity
  conditions than joint density existence.
\item \textbf{Sub-first-order infinities}: extend the indexed
  notation to parameterize logarithmic and iterated-logarithmic
  infinities, which arise in the harmonic series
  (Section~\ref{sec:limitations}) and asymptotic analysis.
\item \textbf{Functorial structure}: investigate whether the
  cross-domain pattern of Section~\ref{sec:crossdomain} admits
  a categorical formalization---specifically, whether the index
  map defines a functor between suitable categories of singular
  expressions and $K^* \times \mathbb{Z}$.
\end{enumerate}

\subsection{The Vision}

Before Bombelli, $\sqrt{-1}$ was impossible. After Gauss, it was $i$.
The notation changed everything---not because new mathematics was
discovered, but because existing mathematics became usable.

Division by zero is at a similar point. The mathematics exists. The
consistency is proven. The computations check out. What remains is
adoption: calculators that output $\vi{5}$ instead of \texttt{ERROR},
textbooks that teach derivatives without limits, and a generation of
students for whom $5/0 = \vi{5}$ is as natural as $\sqrt{-1} = i$.


% ============================================================
% Bibliography
% ============================================================

\begin{thebibliography}{99}

% === Foundational (must-cite) ===

\bibitem{robinson1966}
A.~Robinson, \emph{Non-Standard Analysis}, North-Holland, Amsterdam, 1966.

\bibitem{keisler1976}
H.~J.~Keisler, \emph{Elementary Calculus: An Infinitesimal Approach},
Prindle, Weber \& Schmidt, 1976.
Available free at \url{https://people.math.wisc.edu/~hkeisler/calc.html}.

\bibitem{carlstrom2004}
J.~Carlstr\"om, ``Wheels---On division by zero,''
\emph{Mathematical Structures in Computer Science}, vol.~14, no.~1,
pp.~143--184, 2004.

\bibitem{bergstra2019}
J.~A.~Bergstra, ``Division by zero: A survey of options,''
\emph{Transmathematica}, 2019.

% === Closest prior art ===

\bibitem{santangelo2016}
B.~Santangelo, ``A new algebraic structure that extends fields and allows
for a true division by zero,'' arXiv:1611.06838, 2016.

% === Grossone debate ===

\bibitem{sergeyev2003}
Ya.~D.~Sergeyev, \emph{Arithmetic of Infinity}, Edizioni Orizzonti
Meridionali, Cosenza, 2003.

\bibitem{sergeyev2017}
Ya.~D.~Sergeyev, ``Numerical infinities and infinitesimals: Methodology,
applications, and repercussions on two Hilbert problems,''
\emph{EMS Surveys in Mathematical Sciences}, vol.~4, no.~2, pp.~219--320, 2017.

\bibitem{sergeyev2019}
Ya.~D.~Sergeyev, ``Independence of the grossone-based infinity methodology
from non-standard analysis and its consistency,''
\emph{Foundations of Science}, vol.~24, no.~1, pp.~187--204, 2019.

\bibitem{gutman2017}
A.~E.~Gutman, M.~G.~Katz, T.~S.~Kudryk, and S.~S.~Kutateladze,
``The Mathematical Intelligencer flunks the Olympics,''
\emph{Foundations of Science}, vol.~22, no.~3, pp.~539--555, 2017.

% === Numerosity ===

\bibitem{benci2003}
V.~Benci and M.~Di~Nasso, ``Numerosities of labelled sets: A new way of
counting,'' \emph{Advances in Mathematics}, vol.~173, no.~1, pp.~50--67, 2003.

\bibitem{benci2019}
V.~Benci and M.~Di~Nasso, \emph{How to Measure the Infinite: Mathematics
with Infinite and Infinitesimal Numbers}, World Scientific, 2019.

% === Other frameworks ===

\bibitem{conway1976}
J.~H.~Conway, \emph{On Numbers and Games}, Academic Press, London, 1976.

\bibitem{kock2006}
A.~Kock, \emph{Synthetic Differential Geometry}, 2nd ed.,
Cambridge University Press, 2006.

\bibitem{bell1998}
J.~L.~Bell, \emph{A Primer of Infinitesimal Analysis},
Cambridge University Press, 1998.

% === Alternative division-by-zero approaches ===

\bibitem{saitoh2014}
S.~Saitoh, ``Generalized inversions of Hadamard and tensor products for
matrices,'' \emph{Advances in Linear Algebra \& Matrix Theory}, vol.~4,
no.~2, pp.~87--95, 2014.

\bibitem{bergstra2009}
J.~A.~Bergstra, Y.~Hirshfeld, and J.~V.~Tucker, ``Meadows and the
equational specification of division,'' \emph{Theoretical Computer Science},
vol.~410, no.~12--13, pp.~1261--1271, 2009.

% === Divergent series ===

\bibitem{tao2010}
T.~Tao, ``The Euler--Maclaurin formula, Bernoulli numbers, the zeta function,
and real-variable analytic continuation,'' blog post, April 10, 2010.
Available at \texttt{terrytao.wordpress.com}.

% === Computer arithmetic ===

\bibitem{ieee754}
IEEE, ``IEEE Standard for Floating-Point Arithmetic,''
IEEE Std 754-2019, 2019.

\bibitem{fog2019}
A.~Fog, ``NaN propagation and payloads in IEEE 754 floating-point arithmetic,''
ForwardCom project documentation, 2019.

% === Infinite set sizes ===

\bibitem{wenmackers2024}
S.~Wenmackers, ``Comparing infinities: Six methods and their philosophical
implications,'' \emph{Journal for the Philosophy of Mathematics}, 2024.

\bibitem{feng2026aletheia}
T.~Feng, T.~H.~Trinh, G.~Bingham, D.~Hwang, Y.~Chervonyi, J.~Jung,
J.~Lee, C.~Pagano, S.-h.~Kim, F.~Pasqualotto, S.~Gukov, J.~N.~Lee,
J.~Kim, K.~Hou, G.~Ghiasi, Y.~Tay, Y.~Li, C.~Kuang, Y.~Liu, H.~Lin,
E.~Z.~Liu, N.~Nayakanti, X.~Yang, H.-T.~Cheng, D.~Hassabis,
K.~Kavukcuoglu, Q.~V.~Le, and T.~Luong,
``Towards autonomous mathematics research,''
arXiv:2602.10177, 2026.

% === Phase 4 additions ===

\bibitem{jacobs2021}
B.~Jacobs, ``Multinomial and hypergeometric distributions in Markov
categories,'' in \emph{Proceedings of POPL 2021},
arXiv:2101.03391.

\bibitem{todorov1990}
T.~D.~Todorov, ``An existence result for a class of partial differential
equations with smooth coefficients,'' \emph{Proceedings of the AMS},
vol.~109, no.~1, pp.~1--10, 1990.

\bibitem{vernaeve2025}
H.~Vernaeve, ``Nonstandard analysis and the theory of distributions,''
arXiv:2510.16484, 2025.

\bibitem{albeverio1986}
S.~Albeverio, J.~E.~Fenstad, R.~H\o egh-Krohn, and T.~Lindstr\o m,
\emph{Nonstandard Methods in Stochastic Analysis and Mathematical Physics},
Academic Press, 1986.

\bibitem{meyenburg2025}
J.~Meyenburg, ``On indexed zeros and infinities,''
\emph{International Journal of Mathematics Trends and Technology},
vol.~71, 2025.

\bibitem{fereydoni2025}
M.~Fereydoni, ``Algebraic fixed-point stratification of infinities,''
ResearchGate preprint, 2025.

\end{thebibliography}


% ============================================================
\appendix
\section{Verification Details}
\label{app:verification}
% ============================================================

Every mathematical claim in this paper has been verified by at least
one computational tool. The full verification suite is available in
the supplementary repository at \texttt{verification/}, with detailed
per-tool outputs saved in \texttt{verification/\_results/}. To
reproduce: \texttt{python3 verification/run\_all.py}.

\subsection{Verification Categories}

An internal audit of the original verification suite (April 2026)
revealed that some checks tested classical mathematics rather than
IVNA's axiom system. The suite was rebuilt with honest categorization:

\begin{table}[h]
\centering
\begin{tabular}{llrl}
\toprule
\textbf{Category} & \textbf{Tool} & \textbf{Checks} & \textbf{Result} \\
\midrule
A: IVNA-native       & Python (ivna.py)          & 185 & 185/185 pass \\
A: Core unit tests   & Python (ivna.py)          & 30  & 30/30 pass \\
B: NSA embedding     & SymPy symbolic            & 70  & 70/70 pass \\
C: Classical corresp. & SymPy                    & 93  & 93/93 pass \\
Z3 axiom encoding    & Z3 SMT solver             & 13  & 13/13 pass \\
Wolfram cross-verif. & Wolfram Engine            & 42  & 42/42 pass \\
Formal proofs        & Lean 4.16                 & 23  & Compiles (0 sorry) \\
Meta-verification    & Python                    & 18  & 18/18 pass \\
\midrule
\textbf{Total}       & \textbf{6 independent chains} & \textbf{403+Lean} & \textbf{0 failures} \\
\bottomrule
\end{tabular}
\caption{Verification summary by category. Category~A checks exercise
IVNA's \texttt{Virtual} class directly. Category~B proves axiom
consistency under the NSA embedding. Category~C confirms that IVNA
notation maps correctly onto classical results (honest framing---these
are correspondence checks, not independent derivations).
Meta-verification confirms the suite itself is structurally sound.}
\label{tab:verification}
\end{table}

\subsection{Lean~4 Formalization}

The Lean~4 project (\texttt{lean-ivna/}) contains four files:

\begin{itemize}
\item \texttt{Basic.lean}: The inductive type \texttt{V F} representing
  virtual numbers over a field \texttt{F}, with constructors
  \texttt{real}, \texttt{zero}, and \texttt{inf}.

\item \texttt{Axioms.lean}: The \texttt{IVNASystem F} structure encoding
  all 11 core axioms as equational constraints, plus commutativity
  of multiplication.

\item \texttt{Model.lean}: A concrete instance \texttt{IVNASystem Bool}
  over GF(2), proving that the axiom set is satisfiable (no
  contradiction). This is a machine-checked consistency proof.

\item \texttt{Theorems.lean}: Twelve theorems derived from the axioms
  alone, valid in \emph{any} model. Key results include:
  \begin{itemize}
  \item \textbf{Theorem 1} (Division-by-zero roundtrip):
    $(\text{real}(y) / \vz{x}) \cdot \vz{x} = \text{real}(y)$,
    proven in 10 equational rewriting steps.
  \item \textbf{Theorem 4--5} (Self-division):
    $\vz{x} / \vz{x} = 1$ and $\vi{x} / \vi{x} = 1$.
  \item \textbf{Theorem 6--7} (Scalar associativity):
    $n \cdot (m \cdot \vz{x}) = (nm) \cdot \vz{x}$.
  \item \textbf{Theorem 9--12} (Addition commutativity and associativity).
  \end{itemize}
\end{itemize}

\noindent
To verify: install Lean~4.16+ via \texttt{elan}, then run
\texttt{cd lean-ivna \&\& lake build}. Build completes in approximately
10 seconds with zero errors.

\subsection{Z3 Axiom Encoding}

The Z3 SMT solver encodes IVNA's index arithmetic in first-order
real arithmetic and verifies satisfiability, roundtrip properties,
and axiom independence (13 checks total). Notable findings:

\begin{itemize}
\item The roundtrip $(y / \vz{x}) \cdot \vz{x} = y$ is a tautology
  of real arithmetic under the NSA embedding (negation is UNSAT).
\item The product rule constant must be 1: enforcing
  $c \cdot xy = xy$ with $c \neq 1$ is UNSAT.
\item D-INDEX-ZERO ($\vz{x} - \vz{x} = 0$) is \emph{derivable}
  from Axiom~\ref{ax:zadd} and the embedding ($0 \cdot \varepsilon = 0$),
  and is therefore a convenience axiom rather than an independent postulate.
\item The rejected variant $\vz{1} \cdot \vi{1} = 2\pi$ is UNSAT
  when combined with Axiom~\ref{ax:zi}.
\end{itemize}

\subsection{Python Test Suite}

The Python implementation (\texttt{code/ivna.py}) serves as a reference
implementation and test harness. The 30 core tests cover:

\begin{itemize}
\item All 11 axioms with multiple index values (including negative,
  fractional, and irrational indices)
\item Structural properties: associativity (3 type combinations),
  commutativity, distributivity
\item The division-by-zero roundtrip for 28 distinct $(y, x)$ pairs
\item Higher-order interactions up to order 3
\item The D-INDEX-ZERO rule (index cancellation exits to real~0)
\item Derivatives of $x^n$ for $n = 2, 3, 4, 5$ at multiple evaluation points
\item Transcendental derivatives: $\sin$, $\cos$, $e^x$, $\ln$, $1/x$
  (each exercising the A-VT $\to$ A8 pipeline, not a passthrough)
\item Derivative machinery verification: confirms the derivative
  function calls the Virtual Taylor constructor, divides via
  Axiom~\ref{ax:zzdiv}, and produces higher-order virtual zero
  residuals
\item The exponential axiom: $(1 + \vz{x})^{\vi{y}} = e^{xy}$ for
  6 index pairs
\end{itemize}

\noindent
To run: \texttt{python3 code/ivna.py}. Expected output:
\texttt{30 passed, 0 failed}.

\subsection{SymPy Verification}

Two SymPy-based verification files cover distinct purposes:

\begin{itemize}
\item \textbf{Category B: NSA embedding} (70 checks,
  \texttt{cat\_b\_nsa\_embedding.py}): verifies all axioms A1--A11,
  D-INDEX-ZERO, A-EXP, negation, and higher-order virtual behavior
  (orders 2--4) under the hyperreal mapping $\vz{x} = x\varepsilon_0$,
  $\vi{x} = x/\varepsilon_0$.
\item \textbf{Category C: Classical correspondence} (93 checks):
  confirms that IVNA notation maps correctly onto known results
  across nine domains (Bayes, Borel--Kolmogorov, Dirac delta,
  removable singularities, infinity subtraction, residues,
  compound growth, blow-ups, KL divergence). Honestly labeled
  as correspondence checks---notation, not independent derivation.
\end{itemize}

\subsection{Wolfram Cross-Verification}

The Wolfram Engine verification (42 checks,
\texttt{wolfram\_crosscheck.py}) independently confirms all core
axioms, roundtrip properties, the exponential axiom, five derivative
classes, higher-order embedding behavior, and classical correspondence
results (limits, residues, Dirac delta)---using a different symbolic
engine from SymPy.

\subsection{Meta-Verification}

The meta-verification layer (18 checks, \texttt{meta\_verify.py})
checks the verification suite itself: that Category~A files import
\texttt{ivna.py} (actually test the system), that Category~C files
do not (honest about being classical), that no assertions are
tautological, that Z3 contains no trivial checks, and that
\texttt{ivna\_derivative()} exercises the A-VT pipeline rather than
being a passthrough. This layer runs first---if it fails, the
remaining results are not trusted.

\end{document}
